\documentclass{article}
\usepackage[T2A]{fontenc}
\usepackage[utf8]{inputenc}
\usepackage[russian]{babel}
\usepackage[margin=2cm]{geometry}
\usepackage{icomma}
\usepackage{amsmath}
\usepackage{amssymb}
\usepackage{enumitem}
\usepackage{graphicx}
\usepackage{url}
\usepackage{xfrac}
\usepackage{fancyvrb}

\pagestyle{empty}

\begin{document}

% ==============================================================================

\begin{flushright}
    \textit{Выполнил: Максимов Н. Д.} \\
    \textit{Группа: 09-414} \\
    \textit{Вариант: 8}
\end{flushright}

\begin{center}
    \bfseries\LARGE Отчет по лабораторной работе №3
\end{center}
\bigskip

% ==============================================================================

\paragraph{Постановка задачи} По заданию необходимо построить двусторонний доверительный интервал для дисперсии нормального распределения  при заданной надежности $Q=0.95$. Для построения двустороннего доверительного интервала сначала определяется уровень значимости $\alpha = 1 - Q$. Так как интервал симметричен относительно вероятности попадания в «хвосты» распределения, значение $\alpha$ делится пополам. Таким образом, для расчетов используется величина $\alpha/2$, отсекающая соответствующую область значений критических точек. Верхняя и нижная доверительные границы вычисляеются по формулам:
\[ \overline{\sigma^2} = \frac{n}{\chi_{(\sfrac\alpha2)}}S^2 \qquad\qquad\qquad \underline{\sigma^2} = \frac{n}{\chi_{(1 - \sfrac\alpha2)}} S^2, \]
где $n$ --- число элементов в выборке, $S^2$ ---  ее выборочная дисперсия (смещенная), $\chi$ --- хи-квадрат распределение с $n-1$ степенями свободы.

% ==============================================================================

\paragraph{Полученные значения} В результате проведенных вычислений для заданного уровня значимости $\alpha = 0.05$ и смещенной выборочной дисперсии $S^2 = 17.515$ был получен следующий доверительный интервал:
\[ \sigma^2 \in [12.936; \, 26.007] \]

% ==============================================================================

\paragraph{Листинг}

\noindent \texttt{main.py}
\begin{Verbatim}[tabsize=4, fontsize=\small]
def mean(data: list[float]) -> float:
    return sum(data) / len(data)

def variance(data: list[float], unbiased: bool = False) -> float:
    if unbiased:
        return len(data) / (len(data) - 1) * variance(data, unbiased=False)
    return sum(x**2 for x in data) / len(data) - mean(data)**2

def variance_interval(data: list[float], alpha: float) -> tuple[float, float]:
    n = len(data)
    df = n - 1
    s2 = variance(data, unbiased=False)

    lower = n * s2 / stats.chi2.ppf(1 - alpha / 2, df)
    upper = n * s2 / stats.chi2.ppf(alpha / 2, df)
    return lower, upper
\end{Verbatim}

% ==============================================================================

\vfill
\noindent\rule{\textwidth}{0.4pt}
\small Полный исходный код лабораторной работы доступен в репозитории на GitHub: \url{https://github.com/mndkpfu/MS}

% ==============================================================================

\end{document}